\appendix

\section{\textcolor{blue}{Reversible construction of the point-arithmetic circuit}}
\label{app:oracle}

\begin{revision}
The construction uses the controlled-addition framework of
\cite{ProosZalka2003,Roetteler2017} and the standard elliptic-curve
group law \cite{Silverman2009,Washington2008}. The encoding in
Equation~\eqref{eq:point_encoding}, branch organisation, and cleanup
below specify the reference implementation. For valid point encodings,
\begin{equation}
\label{eq:reversible_circuit}
\begin{aligned}
U_{f_P}:\quad
&\lvert k\rangle\lvert\mathrm{enc}(R)\rangle
 \lvert0\rangle_{\mathrm{aux}}\\
&\longmapsto
 \lvert k\rangle\lvert\mathrm{enc}(R+kP)\rangle
 \lvert0\rangle_{\mathrm{aux}}.
\end{aligned}
\end{equation}
Initializing $R=\mathcal O$ yields Equation~\eqref{eq:oracle}.
\end{revision}

\subsection{\textcolor{blue}{Modular arithmetic blocks}}
\begin{revision}
Let $\ell=\lceil\log_2p\rceil$. Comparisons write predicates into
flags by XOR, preserve their inputs, and reverse temporary subtraction
work. For a classical constant $a\in\{0,\ldots,p-1\}$, modular addition
implements
\begin{equation}
\label{eq:app_r5_constant_addition}
\begin{aligned}
&\lvert x\rangle\lvert0\rangle_{\mathrm{aux}}\\
&\longmapsto\lvert(x+a)\bmod p\rangle
 \lvert0\rangle_{\mathrm{aux}},\quad x<p.
\end{aligned}
\end{equation}
Inputs $x\geq p$ are fixed. On valid inputs, an extended register
holds $x+a$, and $p$ is subtracted when required. For the output
$u=(x+a)\bmod p$, reduction occurred exactly when $u<a$; this
predicate clears the reduction flag. Repeating the range test clears
the validity flag. Addition of a quantum register is assembled by
letting its bit $y_j$ control addition of $2^j\bmod p$, preserving $y$.
Circuit reversal gives subtraction.

Multiplication preserves both operands and uses a modular accumulator:
\begin{equation}
\label{eq:app_r5_multiply_accumulate}
\begin{aligned}
&\lvert x\rangle\lvert y\rangle\lvert t\rangle
 \lvert0\rangle_{\mathrm{aux}}\\
&\longmapsto\lvert x\rangle\lvert y\rangle
 \lvert(t+xy)\bmod p\rangle
 \lvert0\rangle_{\mathrm{aux}},\quad t<p.
\end{aligned}
\end{equation}
Using $xy=\sum_{i,j}x_i y_j2^{i+j}$, each pair of operand bits
controls addition of $2^{i+j}\bmod p$. Invalid accumulator encodings
are fixed. Squaring follows the corresponding expansion of $x^2$.

Inversion is a clean XOR evaluator:
\begin{equation}
\label{eq:app_r5_inverse_evaluator}
\begin{aligned}
&\lvert w\rangle\lvert t\rangle\lvert0\rangle_{\mathrm{aux}}\\
&\longmapsto\lvert w\rangle\lvert t\oplus I_p(w)\rangle
 \lvert0\rangle_{\mathrm{aux}},
\end{aligned}
\end{equation}
where $I_p(w)=w^{-1}\bmod p$ for $0<w<p$, and zero otherwise.
Here $\oplus$ is bitwise XOR. The reference circuit initializes
$(r_0,r_1,c_0,c_1)=(p,w,0,1)$ and uses active Euclidean rounds:
\begin{equation}
\label{eq:app_r5_euclidean_update}
\begin{aligned}
q&=\lfloor r_0/r_1\rfloor,\\
(r_0,r_1)&\longmapsto(r_1,r_0-qr_1),\\
(c_0,c_1)&\longmapsto(c_1,c_0-qc_1\bmod p).
\end{aligned}
\end{equation}
The invariant $r_i\equiv c_iw\pmod p$ yields the inverse in $c_0$
when $(r_0,r_1)=(1,0)$. There are $2\ell+1$ scheduled rounds,
disabled after $r_1=0$. The compact variant reuses one state quartet
while retaining quotient and activation history. After XORing the
valid inverse into the target, reversal clears all state and history.
This is a local quotient-based implementation, not the optimized
inverter of \cite{Roetteler2017}.
\end{revision}

\subsection{\textcolor{blue}{Point translation and scalar evaluation}}
\begin{revision}
For a finite classical constant $A=(x_A,y_A)$ and ordinary input
$R=(x_R,y_R)$ with $x_R\neq x_A$, evaluate
\begin{equation}
\label{eq:affine_translation}
\begin{aligned}
\lambda&=(y_R-y_A)(x_R-x_A)^{-1}\pmod p,\\
x_{R+A}&=\lambda^2-x_R-x_A\pmod p,\\
y_{R+A}&=\lambda(x_R-x_{R+A})-y_R\pmod p.
\end{aligned}
\end{equation}
Shared differences, inverse, and slope are retained until both coordinates have been XORed into their targets, then uncomputed. Mutually exclusive flags handle $R=\mathcal O$, $R=-A$, and $R=A$ with $A\neq-A$, producing the constants $A$, $\mathcal O$, and $2A$. When $A=-A$, the opposite-point branch takes precedence. Thus only the ordinary branch uses the nonzero-denominator formula. Validity checks use canonical coordinates, the curve equation, and the identity flag, rather than a list of points.

Let $H_A$ translate valid point encodings and fix invalid ones. Its
inverse is $H_{-A}$. The total point evaluator implements
\begin{equation}
\label{eq:app_r5_total_point_evaluator}
\begin{aligned}
&\lvert z\rangle\lvert t\rangle\lvert0\rangle_{\mathrm{aux}}\\
&\longmapsto\lvert z\rangle\lvert t\oplus H_A(z)\rangle
 \lvert0\rangle_{\mathrm{aux}}.
\end{aligned}
\end{equation}
An invalid input is therefore copied by XOR into the target. Starting
with a zero target, evaluate $H_A$, exchange registers, and evaluate
$H_{-A}$ to cancel the old encoding. This is
Equation~\eqref{eq:point_update_cleanup}; the zero target is workspace,
not the identity encoding. Controlling every stage gives $V_A$ and
fixes invalid encodings under the completed in-place update. Branch
flags are uncomputed after use, and non-overlapping arithmetic scratch
allocations share a pool.

Precompute $P_0=P$, $P_{j+1}=2P_j$ using at most $m-1$ classical
doublings, without knowing $r$ or enumerating scalar outputs. Initialize
the accumulator at $\mathrm{enc}(\mathcal O)$ and apply $V_{P_j}$
with control $k_j$, omitting $P_j=\mathcal O$. After $j$ bits,
\begin{equation}
\label{eq:partial_scalar_accumulation}
R_j=\sum_{i=0}^{j-1}k_iP_i,\qquad R_0=\mathcal O.
\end{equation}
Hence $R_m=kP$. Each update preserves the scalar bits and cleans its
workspace for reuse. Linearity extends this action to superpositions.
\end{revision}

\subsection{\textcolor{blue}{Code correspondence and validation}}
\begin{revision}
The modular blocks correspond to
\texttt{modular\_add\_constant},
\texttt{modular\_add\_register},
\texttt{modular\_multiply\_accumulate}, and
\texttt{modular\_square\_accumulate}.
The inverter is \texttt{modular\_inverse\_inplace}, using
\texttt{binary\_quotient}; its external interface remains the clean
XOR evaluator despite the in-place internal Euclidean updates.
The point blocks are \texttt{point\_encoding\_check},
\texttt{point\_addition\_cases},
\texttt{ordinary\_point\_sum\_shared}, and
\texttt{point\_sum\_evaluator\_compact}.
The in-place and controlled translations are
\texttt{point\_add\_inplace} and \texttt{controlled\_point\_add},
using \texttt{evaluator\_variant="compact"}.
\texttt{EllipticOrderFinding} assembles the complete order-finding
circuit. Standard reversible gates underlie these local assemblies
\cite{Barenco1995,NielsenChuang2010}.

Tests compare outputs with independent arithmetic references and check
exceptional inputs, invalid encodings, input preservation, and zero
auxiliary outputs. The arithmetic gates are positive-phase basis
permutations, so their verified basis action determines their action
on the tested initialized subspace. Complete noiseless validation
propagates the actual arithmetic gates and derives the Fourier
distribution from their outputs; it is not a hardware experiment.
Recovery and signature statistics are reported separately.
Gate, classical-processing, and workspace bounds are derived in
Appendix~\ref{app:complexity}.
\end{revision}



\section{\textcolor{blue}{Detailed complexity derivations}}
\label{app:complexity}

\subsection{\textcolor{blue}{Parameters and arithmetic costs}}
\begin{revision}
Assume a nonsingular curve over prime $p>3$ and a valid non-identity
point of order $r\geq2$. With $\ell=\lceil\log_2p\rceil$,
Hasse's bound \cite{Hasse1936} gives
\[
r\leq B=p+1+\lfloor2\sqrt p\rfloor.
\]
Algorithm~\ref{alg:hybrid} selects
$m=\lceil\log_2(B^2+1)\rceil$, so $Q=2^m>B^2$ and
$m=\Theta(\ell)$ without assuming that $r$ has $\Theta(\ell)$ bits.

First count each $X$-family instruction, including MCX, as one unit.
A binary constant adder contains at most $\ell(\ell+1)/2$
instructions; comparisons compute and reverse a subtraction.
Modular constant addition therefore costs $O(\ell^2)$.
Register addition uses at most $\ell$ such blocks, and multiplication
at most $\ell^2$. Register comparison and subtraction cost
$O(\ell^3)$; long division uses $O(\ell)$ stages. Each of the
$2\ell+1$ inversion rounds contains quotient and product operations
costing $O(\ell^4)$. Thus, including cleanup,
\begin{center}
\begin{tabular}{lc}
\hline
Reference block & Instructions\\
\hline
Modular constant addition & $O(\ell^2)$\\
Modular register addition & $O(\ell^3)$\\
Product or square accumulation & $O(\ell^4)$\\
Quotient evaluation & $O(\ell^4)$\\
Modular inversion & $O(\ell^5)$\\
Controlled point translation & $O(\ell^5)$\\
\hline
\end{tabular}
\end{center}
The final row uses a constant number of arithmetic blocks, exceptional
branches, and evaluations for in-place cleanup. These are conservative
code-derived bounds, not optimized published resource estimates.
\end{revision}

\subsection{\textcolor{blue}{Quantum gates and workspace}}
\begin{revision}
The final basis is
$\{X,\mathrm{CX},\mathrm{CCX},H,\mathrm{CP}(\theta)\}$,
with ideal controlled phase rotations. An MCX with $k\geq2$ controls
is expanded into $2k-3$ Toffoli gates using $k-2$ clean ancillas;
the partial-conjunction chain is reversed after acting on the target.
A separate ancilla pool is reused between gates. Since the reference
instructions have $O(\ell)$ controls, the expansion gives
\begin{equation}
\label{eq:app_r6_addition_bound}
C_{\mathrm{add}}(\ell)=O(\ell^6).
\end{equation}
At most $m$ controlled additions evaluate $kP$. Preparation and the
inverse QFT then give
\[
G_{\mathrm{one}}=O(m\ell^6+m^2+m+\ell).
\]
Consequently,
\begin{equation}
\label{eq:app_r6_quantum_total}
G_{\mathrm q}=O(T\ell^7),\qquad m=\Theta(\ell).
\end{equation}
This counts the complete evaluation once per execution, not once per
scalar bit. Measurement counts are separate. Rotation synthesis,
depth, routing, and physical error correction are excluded.

The compact inverter retains $O(\ell)$ quotient registers of width
$\ell$ and $O(\ell^2)$ division workspace, despite reusing its
remainder/coefficient quartet. Other arithmetic allocations have no
larger order; the MCX pool adds $O(\ell)$ qubits. Clean workspace is
reused across updates and sequential executions, giving
\begin{equation}
\label{eq:app_r6_workspace}
N_{\mathrm q}=O(m+\ell^2)=O(\ell^2).
\end{equation}
These arguments establish the quantum bounds in
Theorem~\ref{thm:speedup} under the stated ideal gate model.
\end{revision}

\subsection{\textcolor{blue}{Classical processing}}
\begin{revision}
Schoolbook multiplication and division on $\ell$-bit integers cost
$O(\ell^2)$. Allowing $O(\ell)$ Euclidean steps gives conservative
$O(\ell^3)$ bounds for inversion and greatest common divisors.
A classical affine point operation therefore costs $O(\ell^3)$.
Precomputing the constant points takes $O(m)$ point operations,
or $O(m\ell^3)$ bit operations, for valid curve and point inputs.
Offline gate-list construction and synthesis are separate from this
arithmetic precomputation and are not included in the online bounds.

Each measured outcome produces $O(m)$ convergents. Their generation
costs $O(m^3)$ using schoolbook arithmetic. A denominator $v\leq B$
has $O(\ell)$ bits, so checking $vP=\mathcal O$ requires
$O(\ell)$ point operations, or $O(\ell^4)$ bit operations.
Greatest common divisor updates contribute at most $O(m\ell^3)$.
Hence
\begin{equation}
\label{eq:app_r6_postprocessing}
C_{\mathrm{post}}=O(m^3+m\ell^4).
\end{equation}
No subgroup enumeration or trial factorization is performed.

A nonzero estimate $\hat r$ is bounded by $B$. One signing attempt
therefore costs $O(\ell^4)$, dominated by scalar multiplication,
for an $O(\ell)$-bit representative $z$ supplied as input. Hashing
and initial key generation are excluded; nonce retries do not repeat
order recovery. For $m=\Theta(\ell)$,
\begin{equation}
\label{eq:app_r6_classical_total}
\begin{aligned}
C_{\mathrm{hyb,cl}}
&=O\bigl(m\ell^3+T(m^3+m\ell^4)+\ell^4\bigr)\\
&=O(\ell^4+T\ell^5).
\end{aligned}
\end{equation}
Classical bit operations and quantum gate counts remain separate;
this is not a competing classical order-finding algorithm's cost.
\end{revision}

\subsection{\textcolor{blue}{Recovery probability and resulting bounds}}
\begin{revision}
Assume ideal arithmetic, an injective point encoding, an exact inverse
QFT, and independent executions. Following the standard order-finding
argument \cite{Shor1997}, let $y_s$ be nearest to $Qs/r$ for
$\gcd(s,r)=1$. Since $Q>B^2\geq r^2$,
\[
\left|\frac{y_s}{Q}-\frac{s}{r}\right|
\leq\frac{1}{2Q}<\frac{1}{2r^2}.
\]
The convergent criterion \cite{HardyWright2008} therefore yields
the denominator $r$, which passes the group-identity test.

For completeness, a conservative probability bound follows by
partitioning the $Q$ inputs into $r$ residue classes of lengths $N_a$.
Their point outputs are orthogonal and $\sum_aN_a=Q$. At $y_s$,
the geometric-sum step satisfies $|\theta|\leq\pi r/Q$ and
$N_a|\theta|/2\leq3\pi/4$. Thus each sum has magnitude at least
$N_a/4$, using
$\sin(u)/u\geq2\sqrt2/(3\pi)>1/4$ for $|u|\leq3\pi/4$,
with the continuous value at zero.
Consequently,
\begin{equation}
\label{eq:app_r6_success_lower_bound}
\begin{aligned}
\Pr(y_s)&\geq\frac{\sum_aN_a^2}{16Q^2}
\geq\frac{1}{16r},\\
\alpha&\geq\frac{\varphi(r)}{16r},
\end{aligned}
\end{equation}
where $\alpha$ is the probability of including $r$ in one execution
and $\varphi$ is Euler's totient function. The selected outcomes are
distinct. All accepted denominators are multiples of $r$, so including
$r$ fixes their accumulated greatest common divisor at $r$.

The standard estimate
$\varphi(r)/r=\Omega(1/\log\log r)$ used in \cite{Shor1997},
together with $r\leq B=\Theta(p)$, gives
$\alpha\geq c/(1+\log\ell)$ for some $c>0$, adjusting for small
orders. Therefore
\[
\Pr(\hat r\neq r)
\leq e^{-cT/(1+\log\ell)},
\]
including both absent and nonminimal candidates as failures.
For $0<\varepsilon<1/2$, a sufficient budget scales as
\begin{equation}
\label{eq:app_r6_measurement_budget}
T=O\!\left((1+\log\ell)\ln\frac{1}{\varepsilon}\right).
\end{equation}
This is an asymptotic sufficient budget, not an empirical success
rate or a minimality certificate supplied to signing.

For fixed $\varepsilon$, substitution gives
\begin{equation}
\label{eq:app_r6_total_confidence_bounds}
\begin{aligned}
G_{\mathrm q}&=\bigOt(\ell^7),\\
C_{\mathrm{hyb,cl}}&=\bigOt(\ell^5),\\
N_{\mathrm q}&=O(\ell^2).
\end{aligned}
\end{equation}
Here $\bigOt$ hides polylogarithmic factors in $\ell$. These bounds
describe the reference implementation, not optimality, hardware
noise, or cryptographic security. Comparisons with alternative
classical methods require their own cost assumptions.
\end{revision}




\section{Worked instance over \texorpdfstring{$E(\mathbb{F}_{29})$}{E(F29)}}\label{app:instance}

\begin{revision}
For reproducibility, the numerical instance uses $p=29$, $a=0$, $b=7$, and $G=P=(6,7)$. Direct enumeration gives $\#E(\mathbb{F}_{29})=30$, including the identity $\mathcal{O}$. The point $P$ has order $r=5$, since $5P=\mathcal{O}$ and $kP\neq\mathcal{O}$ for $1\leq k<5$. The corresponding cofactor is
$h=\#E(\mathbb{F}_{29})/r=30/5=6$. \end{revision}  \begin{revision}
For $p=29$, the coordinate width is $\ell=5$,
and the point register uses $2\ell+1=11$
qubits, including the identity flag.
The Hasse bound gives $B=40$.
Following Algorithm~\ref{alg:hybrid},
$m=\lceil\log_2(B^2+1)\rceil=11$,
so the scalar register has dimension
$Q=2^{11}=2048$.

The periodic function is evaluated through
controlled additions of the precomputed
points $P_j=2^jP$. All scalar inputs are
retained, and outputs equal to $\mathcal O$
are explicitly encoded. No scalar-output
lookup table or domain compression is used.
\end{revision} \begin{revision}
Noiseless validation propagated the reversible
arithmetic gates generated with Qiskit
\cite{Qiskit2023}. All $2048$ scalar inputs
were checked against independent group
arithmetic, including auxiliary cleanup.
The inverse-QFT measurement distribution
was obtained by FFT from the circuit outputs,
without a dense statevector over the full
arithmetic workspace.

The distribution has its largest peaks at
$y\in\{0,410,819,1229,1638\}$, near the
multiples of $Q/5$. Measurement samples
were drawn using random seed $20261012$.
Across $100$ independent recovery trials,
prefix budgets of $T=1,2,4$ measurements
gave empirical exact-order recovery counts
of $77$, $96$, and $100$, respectively.
No parity correction was used.
\end{revision} \begin{revision}
For the fixed signing test, the message
representative is $z=2$, the private key
is $d=2$, and the nonce is $k=1$.
 \textcolor{blue}{The recovered candidate $\hat r=5$ is accepted exactly by the small-prime check, without sampling Miller--Rabin bases.}
Using the recovered estimate $\hat r=5$,
classical multiplication gives $R=(6,7)$,
and the signature components are
$r_{\sigma}=6\bmod5=1$ and
$s=1^{-1}(2+1\cdot2)\bmod5=4$.
Thus $(r_{\sigma},s)=(1,4)$.
For each of the $100$ recovery trials
with $T=4$, the signature generated
using the recovered estimate matched
the independent classical reference
and passed classical arithmetic
verification.
\end{revision}

\begin{revision}
This instance provides reproducible noiseless
validation of reversible point arithmetic,
probabilistic order recovery, and its
integration with classical signature
generation for the stated test parameters.
\end{revision} \begin{revision}
The resource bounds of
Theorem~\ref{thm:speedup} follow from the
circuit analysis in Appendix~\ref{app:complexity},
rather than from this individual numerical
instance.
\end{revision}
